\section{合卷：霸权为什么能维持，又为什么不能变成帝国}

齐、晋、楚的霸权都依赖联盟、礼仪承认和军事威慑。它们能在某个阶段压住冲突，却难以把其他诸侯变成稳定的行政单位。会盟是一种低成本秩序，代价是秩序必须不断重新谈判。

因此，春秋留下的不是统一帝国，而是一套更尖锐的问题：怎样把封国和卿族掌握的资源重新集中？怎样让军役、土地和行政不只依赖贵族关系？这些问题将在战国各国的变法中得到截然不同的回答。
